% Slides 21--22: conclusions and Q&A.

\begin{frame}{What the Project Establishes}
  \begin{columns}[T,onlytextwidth]
    \begin{column}{0.49\textwidth}
      \keypoint{Reference conclusions}
      \begin{itemize}
        \item Distinct removal rates remove the simple conservation law.
        \item Solutions are nonnegative, bounded, and dissipative; invasion
              thresholds organize washout and coexistence.
        \item Persistence can hold and Hopf bifurcation can create cycles.
      \end{itemize}

      \vspace{0.18cm}
      \keypoint{Our numerical conclusions}
      \par The simulator reproduces $P_0$--$P_3$, damped oscillations, and
      removal-driven regime changes with robust RK4 results.
    \end{column}
    \begin{column}{0.47\textwidth}
      \keypoint{Limitations}
      \begin{itemize}
        \item Parameters are constructed examples, not experimental estimates.
        \item The project illustrates but does not continue a Hopf branch.
        \item Initial-condition tests cannot prove global persistence, and
              calibration data are unavailable.
      \end{itemize}

      \vspace{0.28cm}
      \keypoint{Central lesson}\par
      Every trophic level must grow fast enough on its equilibrium food supply
      to overcome its own removal.
    \end{column}
  \end{columns}
\end{frame}

\begin{frame}[plain]
  \centering
  \vspace{1.25cm}
  {\color{HUSTNavy}\fontsize{40}{46}\selectfont\bfseries Questions?\par}
  \vspace{0.45cm}
  {\color{HUSTTeal}\fontsize{20}{25}\selectfont
    Stability, simulation, sensitivity, or biological interpretation\par}
  \vspace{0.9cm}
  {\fontsize{14}{18}\selectfont\color{MidGray}
    Code and reproducible outputs are available in the project repository.}
\end{frame}
